\section{Conclusion}
  \label{sec:conclusion}

  We have constructed the Lorentzian completion of the spectral entropy framework introduced in previous work.
  While the original formulation was elliptic and sufficient to derive
  the infrared Einstein response and its thermodynamic interpretation, it did not address causal propagation.
  The present paper establishes that the spectral dynamics admits a
  consistent hyperbolic completion with a well-defined retarded Green operator and a real quadratic kernel.

  Linearization about Minkowski spacetime shows that, in the regime
  $R\ell_\chi^2 \ll 1$ and $\omega \ll \ell_\chi^{-1}$,
  the theory propagates exactly two transverse-traceless tensor modes of helicity $\pm 2$.
  Scalar and vector sectors remain non-dynamical in the infrared,
  and the additional pole generated by higher-derivative corrections
  appears at $k^2 \sim \ell_\chi^{-2}$, beyond the frequency range of current gravitational-wave observations.
  The polarization content therefore coincides with that of general relativity in the accessible domain.

  Beyond the leading Einstein sector, the ratio of Seeley--DeWitt
  coefficients fixes a universal $k^4$ correction to the dispersion relation,
  \[
    \omega^2
    =
    c^2 k^2
    -
    \gamma \ell_\chi^2 k^4
    +
    \mathcal{O}(k^6),
  \]
  with $\gamma = 1/(180\zeta)$ and $\zeta=\mathcal{O}(1)$ encoding scheme-dependent normalization factors.
  In the natural spectral scheme, $\gamma=1/30$.
  This structure provides a direct mapping to the parametrizations used
  in gravitational-wave catalogues and yields a bound on the pre-geometric scale $\ell_\chi$ from observational data.

  The Lorentzian completion therefore closes the logical chain of the spectral program:
  the elliptic structure encodes spatial correlations,
  the renormalized first variation yields the infrared Einstein response,
  and the hyperbolic extension ensures causal propagation and observational testability.
  Gravitational waves thus provide a concrete infrared probe of the
  pre-geometric scale underlying the spectral construction.

  Further developments include a fully covariant coupling to matter
  sectors and a systematic analysis of non-local corrections beyond the leading $k^4$ term.
  These directions may sharpen the connection between the spectral
  substrate and precision gravitational observations.
